\documentclass[11pt]{article}
% Shared style for BELAND-PLUS academic revision R3
% Typography: Latin Modern (the canonical scalable Computer Modern family).
\usepackage[T1]{fontenc}
\usepackage[utf8]{inputenc}
\usepackage{lmodern}
\usepackage{microtype}
\usepackage{amsmath,amssymb,mathtools,bm}
\usepackage{booktabs,threeparttable,longtable,tabularx,array,multirow,makecell}
\usepackage{graphicx,adjustbox}
\usepackage{caption,subcaption}
\usepackage{tikz}
\usetikzlibrary{arrows.meta,positioning,fit,calc,backgrounds,matrix,decorations.pathreplacing,shapes.geometric}
\usepackage{pgfplots}
\pgfplotsset{compat=1.18}
\usepackage{xcolor}
\usepackage{geometry}
\geometry{margin=1in,headheight=15pt,footskip=30pt}
\usepackage{setspace}
\setstretch{1.10}
\usepackage{enumitem}
\setlist[itemize]{leftmargin=1.55em,itemsep=0.18em,topsep=0.35em,parsep=0pt}
\setlist[enumerate]{leftmargin=1.8em,itemsep=0.18em,topsep=0.35em,parsep=0pt}
\usepackage{titlesec}
\titleformat{\section}{\large\bfseries}{\thesection}{0.75em}{}
\titleformat{\subsection}{\normalsize\bfseries}{\thesubsection}{0.65em}{}
\titleformat{\subsubsection}{\normalsize\itshape}{\thesubsubsection}{0.55em}{}
\titlespacing*{\section}{0pt}{2.0ex plus .4ex minus .2ex}{0.9ex}
\titlespacing*{\subsection}{0pt}{1.55ex plus .3ex minus .2ex}{0.65ex}
\titlespacing*{\subsubsection}{0pt}{1.1ex plus .2ex minus .2ex}{0.45ex}
\usepackage{fancyhdr}
\pagestyle{fancy}
\fancyhf{}
\fancyhead[L]{\footnotesize\itshape Hurricane Ida and public emergency-service observability}
\fancyhead[R]{\footnotesize BELAND-PLUS}
\fancyfoot[C]{\thepage}
\renewcommand{\headrulewidth}{0.35pt}
\renewcommand{\footrulewidth}{0pt}
\usepackage[round,authoryear]{natbib}
\usepackage{hyperref}
\usepackage[nameinlink,noabbrev]{cleveref}
\usepackage{xurl}
\urlstyle{same}
\usepackage{csquotes}
\usepackage{listings}
\usepackage{tcolorbox}
\tcbuselibrary{breakable,skins}
\usepackage{float}
\usepackage[section]{placeins}
\usepackage{siunitx}
\sisetup{detect-all,group-separator={,},group-minimum-digits=4,table-number-alignment=center}
\usepackage{pdflscape}
\usepackage{etoolbox}

\definecolor{BelandBlue}{HTML}{183A5A}
\definecolor{BelandBlueLight}{HTML}{EAF0F5}
\definecolor{BelandOrange}{HTML}{B5652A}
\definecolor{BelandRed}{HTML}{9E3D36}
\definecolor{BelandGray}{HTML}{4B5563}
\definecolor{BelandLightGray}{HTML}{F4F5F7}
\definecolor{BelandGreen}{HTML}{2E6F64}
\definecolor{BelandGold}{HTML}{A47A24}

\hypersetup{
  colorlinks=true,
  linkcolor=BelandBlue,
  citecolor=BelandBlue,
  urlcolor=BelandBlue,
  pdfborder={0 0 0}
}

\newcommand{\Jzerozero}{J_{00}}
\newcommand{\Jzeroone}{J_{01}}
\newcommand{\Jonezero}{J_{10}}
\newcommand{\Joneone}{J_{11}}
\newcommand{\E}{\mathrm{E}}
\newcommand{\R}{\mathrm{R}}
\newcommand{\Prob}{\Pr}
\newcommand{\Ida}{\mathrm{Ida}}
\newcommand{\Iset}{\mathcal{I}}
\newcommand{\Mset}{\mathfrak{M}}
\newcommand{\Xset}{\mathcal{X}}
\newcommand{\supp}{\operatorname{supp}}
\newcommand{\diam}{\operatorname{diam}}
\newcommand{\argmin}{\operatorname*{arg\,min}}
\newcommand{\argmax}{\operatorname*{arg\,max}}
\newcommand{\ind}{\mathbf{1}}
\newcommand{\Rset}{\mathbb{R}}
\newcommand{\Nset}{\mathbb{N}}
\newcommand{\abs}[1]{\left|#1\right|}
\newcommand{\norm}[1]{\left\lVert#1\right\rVert}
\newcommand{\given}{\,|\,}

\newtcolorbox{plainbox}[1][]{
  enhanced,breakable,colback=BelandBlueLight,colframe=BelandBlue,
  boxrule=0.55pt,arc=1.5mm,left=2.7mm,right=2.7mm,top=1.8mm,bottom=1.8mm,#1}
\newtcolorbox{resultbox}[1][]{
  enhanced,breakable,colback=BelandLightGray,colframe=BelandGray,
  boxrule=0.5pt,arc=1.5mm,left=2.7mm,right=2.7mm,top=1.8mm,bottom=1.8mm,#1}
\newtcolorbox{warningbox}[1][]{
  enhanced,breakable,colback=white,colframe=BelandOrange,
  boxrule=0.65pt,arc=1.5mm,left=2.7mm,right=2.7mm,top=1.8mm,bottom=1.8mm,#1}
\newtcolorbox{definitionbox}[1][]{
  enhanced,breakable,colback=white,colframe=BelandBlue!70,
  boxrule=0.45pt,arc=1.2mm,left=2.5mm,right=2.5mm,top=1.5mm,bottom=1.5mm,#1}

\captionsetup{font=small,labelfont=bf,labelsep=period,skip=5pt,justification=raggedright,singlelinecheck=false}
\renewcommand{\arraystretch}{1.18}
\setlength{\tabcolsep}{5pt}
\setlength{\parindent}{1.35em}
\setlength{\parskip}{0pt}
\setlength{\emergencystretch}{2em}
\raggedbottom

\lstdefinestyle{belandcode}{
  basicstyle=\ttfamily\footnotesize,
  backgroundcolor=\color{BelandLightGray},
  frame=single,
  rulecolor=\color{BelandGray!45},
  breaklines=true,
  breakatwhitespace=false,
  columns=fullflexible,
  keepspaces=true,
  showstringspaces=false,
  xleftmargin=0.5em,
  xrightmargin=0.5em,
  aboveskip=0.7em,
  belowskip=0.7em
}
\lstset{style=belandcode}

\newcommand{\notclaim}[1]{\textit{This does not identify #1.}}
\newcommand{\pp}{percentage points}
\newcommand{\keywords}[1]{\par\noindent\textbf{Keywords:} #1}
\newcommand{\jel}[1]{\par\noindent\textbf{JEL codes:} #1}

\AtBeginEnvironment{tabular}{\small}
\AtBeginEnvironment{tabularx}{\small}
\AtBeginEnvironment{longtable}{\small}

\usepgfplotslibrary{dateplot}
\fancyhead[L]{\footnotesize\itshape Public CAD is not operational ground truth}

\title{\vspace{-1.7em}\textbf{Public CAD Is Not Operational Ground Truth}\\[0.35em]
\large A Hurricane Ida Stress Test of Public Dispatch Data}
\author{Yuwen Zhu\\Carleton University}
\date{\textit{Draft for discussion}\\[0.2em]26 August 2026}

\begin{document}
\maketitle
\thispagestyle{plain}
\vspace{-1.3em}

\begin{plainbox}[title=Abstract]
Researchers commonly measure police response time with timestamps recorded in computer-aided dispatch (CAD) systems. When CAD records are publicly released, that use requires stable field definitions and coverage. In New Orleans, the released record changed abruptly on 28 July 2021. Before that date, none of 419,840 non-officer-initiated records contained an arrival timestamp without a dispatch timestamp. Within two days, that configuration became common, while dispatch-field coverage in officer-initiated records fell from 100 percent to 1 percent. Five weeks later, Hurricane Ida produced a temporary additional shift: among records with a valid arrival field, the public missing-dispatch share reached 43.4 percent on 31 August and 49.3 percent on 1 September. The largest half-day change was 53.2 percentage points, larger than all 151 post-change ordinary-week comparisons. These results document instability in the released measurement product, not physical police response, police performance, or the mechanism generating the public fields. Public CAD timestamps should be used as operational clocks only after field meaning, coverage, and provenance are validated.
\end{plainbox}

\keywords{calls for service; police responsiveness; administrative data; Hurricane Ida; measurement error; administrative data quality}
\jel{C18, C55, H12, K42}

\section{Introduction}

Researchers and public agencies commonly construct police response-time measures from timestamps recorded in CAD systems. When CAD records are publicly released, that use assumes the fields retain stable meanings and coverage. This paper tests that assumption in New Orleans before treating the released fields as performance measures. It observes whether public dispatch and arrival fields are present and how their joint pattern changes. A missing dispatch timestamp may reflect a missing operational stage, but it may also reflect workflow, later entry, retention, redaction, or export logic. These field-presence measures describe the released record; they are not response times, physical-event indicators, or measures of service quality.

Two empirical facts organize the paper. First, the New Orleans public file changes abruptly on 28 July 2021. Before that date, no non-officer-initiated record has a valid arrival field without a dispatch field. Beginning on 28 July, that configuration appears and persists. At the same time, dispatch fields almost disappear from officer-initiated records even though their arrival fields remain nearly complete. Second, Hurricane Ida produces an additional, temporary reconfiguration within this new regime. On 31 August and 1 September, the public missing-dispatch share---the share missing a dispatch field among records with a valid arrival field---approaches one-half.

The Ida pattern is large relative to later ordinary weeks. Using every non-officer record in each half-day, the largest event-minus-prior-week change in a public field state is 53.2 percentage points. No post-change ordinary-week comparison is as large. A conditional multinomial bootstrap leaves Ida first in all 4,000 draws, although that exercise does not model time-series dependence. A composition-standardized comparison also ranks Ida first, but it retains incomplete common support and is therefore treated as secondary evidence.

The paper's contribution is about measurement, not about the causal effect of Ida. Public records show that the released measurement product changes before the hurricane and changes again during the event. They do not reveal whether the underlying cause was an operational workaround, an entry convention, a retention rule, or an export transformation. The implication is straightforward: a public CAD timestamp should not be treated automatically as an operational clock. Analysts must first establish the relevant denominator, field meaning, coverage, and lineage.

\section{Related literature and contribution}

A large economics and policing literature shows that response delays matter. Faster police arrival can improve crime clearance and reduce injuries, while traffic congestion, staffing, and queueing can slow first responders \citep{BlanesKirchmaier2018,DeAngeloTogerWeisburd2023,BrentBeland2020,MourtgosAdamsNix2024}. These studies demonstrate the value of a valid response-time measure. The present paper asks the logically prior question: whether a public administrative export supplies a stable response clock in the first place.

The paper also relates to work on administrative measurement. Calls-for-service and police records are shaped by reporting, classification, discretion, verification, and data-processing rules \citep{KlingerBridges1997,BoivinCordeau2011,BoivinOuellet2014,SimpsonOrosco2021}. More generally, measurement error and misclassification can distort group comparisons and structural interpretation \citep{KapteynYpma2007,ChenHongNekipelov2011,KnoxLoweMummolo2020}. The evidence here adds a time dimension: the joint presence of public CAD fields changes abruptly, so the mapping from activity to released data is not stable across the sample period.

Disaster and domestic-violence research provides substantive motivation. Hurricane exposure, displacement, psychological stress, partner conflict, and demand for services can move together \citep{AnastarioShehabLawry2009,SchumacherEtAl2010,HarvilleTaylorTesfai2011,FirstEtAl2022}. Police composition can affect reporting \citep{MillerSegal2019}, and pandemic-era studies show that calls, reports, and underlying victimization need not move one-for-one \citep{LeslieWilson2020,BullingerCarrPackham2021,MillerSegalSpencer2022,MillerSegalSpencer2024}. Those studies do not validate a particular CAD timestamp. Ida is used here as a stress test of administrative observability, not as a treatment whose effect on domestic-violence incidence is estimated.

Finally, partial-identification methods provide the appropriate language when public records do not reveal a unique latent process \citep{Manski1989,Manski1990,Manski2003,Tamer2010,Molinari2020}. The paper does not propose a new general partial-identification estimator. It uses that discipline to distinguish quantities directly observed in the public file from quantities that are selected, bounded, or unavailable without additional evidence.

Relative to these literatures, the paper makes three contributions. It documents a date-localized change in the public CAD measurement regime; it shows that Ida produces an additional extreme but temporary reconfiguration within the later regime; and it translates those facts into practical requirements for measuring response time and reported-DV performance.

\section{Data and public field measures}

\subsection{Official files and analysis population}

The analysis uses the official New Orleans Calls for Service files for 2020--2026 \citep{DataNOLA2020,DataNOLA2021,DataNOLA2022,DataNOLA2023,DataNOLA2024,DataNOLA2025,DataNOLA2026}. The reproduction package binds the 2020--2024 annual files to official dataset identifiers, row counts, observed date ranges, monthly source hashes, and deterministic annual bundle hashes.

The main analysis is restricted to records coded \texttt{SelfInitiated=N}, which the public data distinguish from items generated by officers in the field. These are referred to as \emph{non-officer-initiated} records. Officer-initiated records are analyzed separately because their public dispatch-field convention changes sharply in July 2021.

For each record, define the dispatch field as present when \texttt{TimeDispatch} is nonblank. Define the arrival field as valid when \texttt{TimeArrive} parses and is no earlier than \texttt{TimeCreate}. These two indicators produce four public configurations.

\begin{table}[!htbp]
\centering
\caption{Four configurations in the released public record}
\label{tab:states}
\begin{tabularx}{\textwidth}{@{}lccX@{}}
\toprule
Public configuration & Dispatch field & Valid arrival field & What the label does not establish \\
\midrule
Neither field & No & No & No call activity or officer action. \\
Dispatch only & Yes & No & No physical arrival. \\
Arrival only & No & Yes & Physical arrival without dispatch. \\
Both fields & Yes & Yes & Stable endpoint meanings or a complete operational pathway. \\
\bottomrule
\end{tabularx}
\begin{minipage}{0.96\textwidth}\footnotesize
These labels describe the released record. ``Arrival only'' does not mean that an officer physically arrived without being dispatched.
\end{minipage}
\end{table}

Three descriptive series summarize the public file: the share of records with a dispatch field, the share with a valid arrival field, and the \emph{public missing-dispatch share}---among records with a valid arrival field, the share whose public dispatch field is missing. Calls are grouped by creation date; dispatch and arrival fields are read from the eventual released record.

\section{A public-data regime change on 28 July 2021}

The arrival-only configuration is completely absent from non-officer-initiated records before 28 July 2021. Across 419,840 records created from January 2020 through 27 July 2021, its count is exactly zero. It appears abruptly on 28 July and remains present thereafter.

A direct audit of the official July 2021 compressed CSV confirms that this is not an artifact of the aggregate pipeline. All nonblank dispatch and arrival fields in the audit interval parse under the stated rule, and the daily state counts agree exactly with the packaged aggregate tally. Table~\ref{tab:julybreak} shows the central transition.

\begin{table}[!htbp]
\centering
\caption{The public-file transition, 27--29 July 2021}
\label{tab:julybreak}
\begin{tabular}{@{}lrrr@{}}
\toprule
 & 27 July & 28 July & 29 July \\
\midrule
Non-officer records & 731 & 620 & 666 \\
Non-officer arrival-only records & 0 (0.0\%) & 34 (5.5\%) & 69 (10.4\%) \\
Non-officer dispatch-field share & 90.0\% & 82.1\% & 76.6\% \\
Officer records & 479 & 575 & 504 \\
Officer dispatch-field share & 100.0\% & 40.0\% & 1.0\% \\
Officer valid-arrival share & 100.0\% & 100.0\% & 100.0\% \\
\bottomrule
\end{tabular}
\begin{minipage}{0.96\textwidth}\footnotesize
Percentages are calculated from the official daily source rows. Full state counts for 25--31 July are reported in the empirical supplement.
\end{minipage}
\end{table}

The simultaneous movement across the two initiation streams is important. Officer-initiated records carry the public initiation label for items generated by officers in the field, yet their public dispatch fields almost disappear while their arrival fields remain complete. The missing public dispatch field therefore cannot be read mechanically as the absence of an officer or of all operational activity. For this stream, the post-break change is necessarily a change in production of the released record, although the public evidence does not identify the specific entry, retention, transformation, or export step. The discrepancy lies somewhere between activity and released representation.

Figure~\ref{fig:monthly} places the break in the longer series. Dispatch-field coverage among non-officer records is about 90 percent in 2020, then declines after July 2021 and reaches 65.4 percent in 2024. The arrival-only configuration, previously absent, becomes increasingly common.

\begin{figure}[!htbp]
\centering
\begin{tikzpicture}
\begin{axis}[
  width=0.94\textwidth,height=6.2cm,
  date coordinates in=x,
  xmin=2020-01-01,xmax=2025-01-01,
  ymin=0,ymax=1,
  xtick={2020-01-01,2021-01-01,2022-01-01,2023-01-01,2024-01-01,2025-01-01},
  xticklabel=\year,
  ylabel={Share among non-officer records},
  grid=major,grid style={BelandLightGray},
  legend style={draw=none,font=\small,at={(0.5,-0.19)},anchor=north,legend columns=3},
  tick label style={font=\small},label style={font=\small}
]
\addplot[BelandBlue,thick] table[x=month,y=dispatch_share,col sep=comma] {r15_monthly_field_completeness.csv};
\addplot[BelandGreen,thick] table[x=month,y=arrival_share,col sep=comma] {r15_monthly_field_completeness.csv};
\addplot[BelandOrange,thick] table[x=month,y=j01_given_arrival,col sep=comma] {r15_monthly_field_completeness.csv};
\draw[BelandRed,dashed,thick] (axis cs:2021-07-28,0) -- (axis cs:2021-07-28,1);
\legend{Dispatch field present,Valid arrival field,Dispatch missing among arrival-observed}
\end{axis}
\end{tikzpicture}
\caption{Public-field completeness among non-officer-initiated records. The dashed line marks 28 July 2021, the first date with an arrival-only record.}
\label{fig:monthly}
\end{figure}

The later annual files confirm that this is a persistent measurement regime rather than a short Ida artifact. In the comparable non-officer denominator, dispatch-field coverage is 65.4 percent in 2024, 66.0 percent in 2025, and 66.8 percent in the 2026 snapshot. The much lower all-row rate of about 41.8 percent combines non-officer records with officer-initiated records, whose dispatch fields are almost always absent under the later convention.

\begin{table}[!htbp]
\centering
\caption{Dispatch-field coverage in the later public-data regime}
\label{tab:current}
\begin{tabular}{@{}lrrr@{}}
\toprule
Year & Non-officer records & Officer records & All public rows \\
\midrule
2024 & 65.4\% & 1.1\% & 42.7\% \\
2025 & 66.0\% & 0.9\% & 41.8\% \\
2026 snapshot & 66.8\% & 0.9\% & 41.9\% \\
\bottomrule
\end{tabular}
\begin{minipage}{0.96\textwidth}\footnotesize
The all-row rate is not comparable to the non-officer series because it pools two streams with different recording conventions.
\end{minipage}
\end{table}

\section{Hurricane Ida within the new regime}

Hurricane Ida made landfall on 29 August 2021, five weeks after the public-file transition. Figure~\ref{fig:idapath} shows the daily field pattern for non-officer records from one week before Ida through the following two weeks.

\begin{figure}[!htbp]
\centering
\begin{tikzpicture}
\begin{axis}[
  width=0.94\textwidth,height=6.1cm,
  date coordinates in=x,
  xmin=2021-08-22,xmax=2021-09-12,
  ymin=0,ymax=1,
  xtick={2021-08-22,2021-08-29,2021-09-03,2021-09-05,2021-09-12},
  xticklabel=\month\ \day,
  ylabel={Share among non-officer records},
  grid=major,grid style={BelandLightGray},
  legend style={draw=none,font=\small,at={(0.5,-0.20)},anchor=north,legend columns=3},
  tick label style={font=\small},label style={font=\small}
]
\path[fill=BelandOrange!12] (axis cs:2021-08-29,0) rectangle (axis cs:2021-09-03,1);
\addplot[BelandBlue,thick,mark=*] table[x=date,y=dispatch_share,col sep=comma] {r15_1_ida_time_path.csv};
\addplot[BelandGreen,thick,mark=square*] table[x=date,y=arrival_share,col sep=comma] {r15_1_ida_time_path.csv};
\addplot[BelandOrange,thick,mark=triangle*] table[x=date,y=j01_given_arrival,col sep=comma] {r15_1_ida_time_path.csv};
\legend{Dispatch field present,Valid arrival field,Dispatch missing among arrival-observed}
\end{axis}
\end{tikzpicture}
\caption{The public-field path around Hurricane Ida. Shading marks the five-day event window. The figure describes the released record, not physical dispatch or arrival.}
\label{fig:idapath}
\end{figure}

During the preceding week, the public missing-dispatch share ranges from 5.1 to 11.7 percent. On landfall day, 29 August, the valid-arrival share falls to 43.8 percent from a preceding-week range of 74--84 percent, while dispatch-field coverage remains 72.9 percent. The public missing-dispatch share then rises to 22.9 percent on 30 August, 43.4 percent on 31 August, and 49.3 percent on 1 September as dispatch-field coverage falls to 47.3 and 40.9 percent on the latter two days. It falls to 16.4 percent on 2 September, returns to its pre-event range by 5 September, and shows one further one-day excursion to 34.7 percent on 10 September before settling. Thus the first event day is marked mainly by missing arrival fields, while the third and fourth days are marked by missing dispatch fields among arrival-observed records.

These movements are large, but their interpretation must remain administrative. The figures index calls by creation date and read fields from the eventual released record. They do not show a live system snapshot, nor do they establish that a unit physically arrived without dispatch. They show that records created during the event were ultimately released in a different field configuration.

\section{How unusual was Ida?}

\subsection{Primary full-count comparison}

For each five-day window, the analysis compares ten twelve-hour cells with the same half-days one week earlier. Within each cell it computes changes in three of the four public configurations: dispatch only, arrival only, and both fields present. The neither-field share is determined by the other three because the four shares sum to one. The summary statistic is the largest absolute change across the thirty cell-by-state comparisons:
\begin{equation}
M(w)=\max_{b,s}\left|\widehat p_{w,b,s}-\widehat p_{w-7,b,s}\right|.
\end{equation}
For Ida, the largest public-state share changed by 53.2 percentage points ($M=0.532$).

The primary reference comparison uses 151 ordinary windows for which both the event and baseline periods occur after 28 July 2021. Prespecified emergency and holiday windows remain excluded. This post-change set was defined after the July transition was identified, so it is a transparent descriptive sensitivity rather than a preregistered test.

The largest reference value is 0.310. Thus Ida is larger than all 151 post-change ordinary-week comparisons. A conditional window-wise multinomial bootstrap independently perturbs Ida and the reference windows. Ida ranks first in all 4,000 draws, with a 95-percent interval of $[0.475,0.601]$ for its statistic. Because the procedure does not preserve serial or cross-window dependence, it is best read as a robustness check against ordinary categorical-count variation, not as formal time-series inference.

\subsection{Secondary standardized comparison}

A second estimator standardizes across common call-type, hour, and schema strata. It also ranks Ida first: the maximum standardized contrast is 0.507, compared with 0.328 for the largest of 153 prespecified references beginning on or after 1 July 2021. The supplement calls this the stage-era set and documents its inherited cutoff. However, the standardized analysis covers only 86.1 percent of event records and 77.0 percent of baseline records. Ida would fail the 0.90 support rule applied to candidate reference windows. The standardized result is therefore secondary; the full-count comparison is the main evidence.

\begin{table}[!htbp]
\centering
\caption{Ida relative to ordinary-week comparisons}
\label{tab:comparison}
\begin{tabularx}{\textwidth}{@{}lrrrX@{}}
\toprule
Design & References & Ida largest change & Largest reference & Interpretation \\
\midrule
Full-count, post-change & 151 & 53.2 pp & 31.0 pp &
Primary descriptive comparison; uses every non-officer record in each half-day. \\
Standardized, stage-era & 153 & 50.7 pp & 32.8 pp &
Same qualitative rank, but incomplete common support makes it secondary. \\
\bottomrule
\end{tabularx}
\end{table}

Two excluded episodes help interpret the result. The week after Ida has a large statistic of 0.478 because its baseline is Ida itself; it captures recovery rather than an independent shock. Hurricane Francine, a later-regime hurricane in September 2024, has a much smaller statistic of 0.110. Two pre-break windows later labelled as Tropical Storm Cristobal and its following comparison week have changes of 0.333 and 0.340, but those movements occur through the dispatch-only and both-field configurations rather than the arrival-only configuration. Ida is therefore not merely a generic consequence of comparing any hurricane week with the preceding week.

All other prespecified statistics and threshold variants are reported in the empirical supplement. They do not all rank Ida first. The main conclusion is narrow: Ida produces an unusually large change in the public field configuration.

\section{What do the data suggest---and what remains unknown?}

\subsection{A recording-convention clue}

The officer-initiated stream provides the most concrete clue. The public definition describes these items as generated by officers in the field. Before 28 July 2021, their dispatch and arrival fields are almost universally present. After the break, arrival remains nearly universal while dispatch falls to approximately 1 percent. The arrival-only configuration is therefore a normal released form for records carrying that public initiation label under the later regime.

During Ida, non-officer records move toward the same public configuration. One plausible pathway is that some calls normally represented with a dispatch field were entered, retained, or exported through a workflow that omitted it. That possibility is consistent with a change in how activity is represented in the released record, but it is not identified. Internal call-origin provenance, event-entry histories, unit-assignment histories, and export logs would be required to distinguish reclassification, operational workarounds, back-filling, and export omission.

\subsection{Disposition mix is not the main accounting explanation}

The late Ida movement also occurs within recorded final-disposition categories. On the afternoon and evening of 31 August, the public dispatch-field share among arrival-observed records falls by 49.5 percentage points. A change in the observed disposition mix would have increased the share by about 1.2 points; within-disposition changes account for approximately $-50.7$ points. The following morning produces nearly the same decomposition. The observed final-disposition mix therefore does not explain the aggregate decline.

This is an accounting result, not a mechanism test. The analysis conditions on records with arrival fields and uses final recorded dispositions. Selection into that population, changes within broad categories, and changes in recording practice remain possible.

\subsection{What public documents establish}

Official metadata define the public field labels and identify the data provider. The DataNOLA page also warns that records are preliminary, may contain mechanical or human error, may change after investigation, and should not be used mechanically for comparison over time \citep{DataNOLA2021}. Public oversight materials document the Ida emergency context \citep{OIPM2021}, while NOPD policy materials describe intended communications and records practices \citep{NOPDPolicies2023}.

The public chronology rules out two proposed technology explanations. The publicly indexed Carbyne APEX cutover occurred in June 2022, eleven months after the July 2021 change \citep{OPCDCarbyne2022}. The Hexagon OnCall Records project was contracted but never launched and therefore cannot explain the transition \citep{NOLAOIGHexagon2024}. These exclusions are useful, but they do not identify the actual cause.

No public versioned changelog or internal-to-public reconciliation was available for this analysis, and no agency confirmation of the July change was obtained. A focused public-records request to OPCD for the CAD export configuration and change history covering July 2021 is the most direct next step for resolving the mechanism. Until such evidence is available, the strongest supported conclusion is a discontinuity in the released record, not a named institutional mechanism.

\section{Implications for response-time and reported-DV research}

The paper's measurement lesson applies beyond Hurricane Ida. Public CAD can support a response-time or reported-DV measure only when the target population and the public endpoints are explicit.

\begin{table}[!htbp]
\centering
\caption{What the public file supports}
\label{tab:implications}
\begin{tabularx}{\textwidth}{@{}p{0.22\textwidth}p{0.34\textwidth}X@{}}
\toprule
Target & What the public file directly supports & Additional evidence needed \\
\midrule
Field completeness &
Shares of records with each released public field, on a declared denominator. &
None beyond source and parsing documentation. \\
Response time &
A duration among records with both declared endpoints. This is a selected-record statistic when endpoints are missing. &
Validated endpoint meanings, coverage on the full denominator, and evidence that the fields map to the intended operational events. \\
Queueing or capacity &
No direct measure from endpoint presence alone. &
Calls holding, queue flows, unit availability, assignments, and internal event histories. \\
Reported-DV performance &
Counts under a declared classification rule and public-field completeness within that classified population. &
Validated classification, treatment of unresolved records, endpoint coverage, and lineage to reports or follow-up stages. \\
\bottomrule
\end{tabularx}
\end{table}

Reported CAD calls are not equivalent to underlying domestic-violence incidence or complete help-seeking. A future DV analysis should preserve unresolved classifications rather than force every record into or out of the target group, and it should report endpoint coverage before interpreting observed durations. Ida motivates this measurement work, but the current system-level analysis does not estimate a DV effect.

\section{Limitations}

The analysis uses eventual public records rather than internal event histories. The July change is verified in the official source bytes, but its institutional cause is not observed. The post-change comparison set was defined after the break was identified and excludes specified emergency and holiday contexts. Reference windows overlap; retaining every second window removes adjacent overlap and leaves Ida first in both alternating phases (supplement). The bootstrap does not preserve all time-series dependence. Public \texttt{SelfInitiated=N} is not equivalent to ``911 only.'' Weather, holidays, mobility, evacuation, reporting behavior, call composition, suppression, and other contemporaneous conditions can affect the observed record.

These limitations restrict interpretation but do not erase the descriptive result. The released field configuration changes abruptly in July 2021, changes further during Ida, and follows different conventions across initiation streams.

\section{Conclusion}

Public CAD files are valuable administrative records, but their timestamps are not operational ground truth by definition. In New Orleans, a new arrival-only configuration appears suddenly on 28 July 2021, accompanied by an almost complete disappearance of dispatch fields from officer-initiated records. Five weeks later, Hurricane Ida produces an additional two-day reconfiguration that is larger than every post-change ordinary-week comparison.

The public evidence identifies the timing and magnitude of changes in the released record. It does not identify the underlying operational sequence or the institutional mechanism that produced those fields. Analysts should therefore treat public CAD timestamps as administrative measures until their meanings, coverage, and lineage have been validated for the specific performance question.

\section*{Data and code availability}

Replication materials, source bindings, the empirical supplement, and the release verifier are available in the public repository at \url{https://github.com/yuwen3756-hue/public-cad-measurement-police-responsiveness-hurricane-ida}. The repository provides derived aggregate artifacts and reproducible code; official raw CAD files remain identified by first-hand public-source links and hashes in the reproduction materials.

\bibliographystyle{apalike}
\bibliography{references_r16_1}

\end{document}
